% Requires amsmath and amssymb in the main report.
% Replace the existing Contract value metrics subsection with this file.
% Equation and subsection numbering are inherited from the main report.
\subsection{Contract value metrics}
\label{sec:q1-contract-value-metrics}

The contract is evaluated from two perspectives: the community's total
cost and the import-limitation service delivered to the DSO. To compare
the three models, define the hourly import quantity
\begin{equation}
G_t^m =
\begin{cases}
\displaystyle\sum_{i\in\mathcal N} I_{i,t}^m,
    & m\in\{\mathcal R,\mathcal I\},\\[4pt]
P_t^{\mathrm{imp},\mathcal C},
    & m=\mathcal C,
\end{cases}
\qquad t\in\mathcal T.
\label{eq:q1-metrics-import}
\end{equation}
Thus, imports are summed across individual meters in $\mathcal R$ and
$\mathcal I$, whereas $\mathcal C$ uses the community's single external
meter. Exports from one member do not offset another member's metered
imports in $\mathcal R$ or $\mathcal I$. Following the one-hour convention,
$G_t^m$ is measured in kWh per interval and is numerically equal to the
average import power in kW over that interval.
For the two contract models the metered totals are
$I_{i,t}^{\mathcal I}=I_{i,t}^{\mathrm{cap},\mathcal I}+x_{i,t}^{\mathcal I}$
and
$P_t^{\mathrm{imp},\mathcal C}=P_t^{\mathrm{imp,cap},\mathcal C}
+x_t^{\mathcal C}$. These totals enter the import metrics and the
energy and tariff charges.

\paragraph{Community perspective.}
The first metric is the optimal total cost $J^m$ (DKK), already defined
by each model's objective. It includes net energy trading costs, grid
tariffs, exceedance penalties where applicable, PV production costs and
load-deviation disutility. Including disutility accounts for the cost
of changing members' preferred consumption schedules.

The benefit of signing a contract relative to the no-contract reference is
\begin{equation}
\Delta J^{m\mid\mathcal R}
= J^{\mathcal R}-J^m,
\qquad m\in\{\mathcal I,\mathcal C\}.
\label{eq:q1-metrics-gain-reference}
\end{equation}
A positive value means that the contract lowers the community's total
cost; a negative value means that remaining without a contract is cheaper.
The additional benefit of signing collectively rather than individually is
\begin{equation}
\Delta J^{\mathcal C\mid\mathcal I}
= J^{\mathcal I}-J^{\mathcal C}.
\label{eq:q1-metrics-gain-pooling}
\end{equation}
This difference measures the combined value of shared contractual capacity,
coordinated dispatch and energy exchanges behind the community meter,
under the same contract terms. These are aggregate benefits; they do not
establish that every member benefits after settlement.

\paragraph{DSO perspective.}
The remaining import above the common cap is measured by
\begin{equation}
V^m
= \sum_{t\in\mathcal T}
  \left[G_t^m-\overline P_t^{\mathrm{DSO}}\right]_+,
\qquad m\in\{\mathcal R,\mathcal I,\mathcal C\},
\label{eq:q1-metrics-exceedance}
\end{equation}
where $[z]_+=\max\{z,0\}$. The unit is kWh under the one-hour convention.
The cap is only an evaluation threshold in $\mathcal R$; it neither
constrains dispatch nor generates a contractual penalty in that model.

Two complementary metrics describe the magnitude and frequency of high
imports:
\begin{align}
P_{\mathrm{peak}}^m
&= \max_{t\in\mathcal T}G_t^m,
\label{eq:q1-metrics-peak}\\
H^m
&= \sum_{t\in\mathcal T}
\mathbf{1}\!\left\{
G_t^m>\overline P_t^{\mathrm{DSO}}+\varepsilon
\right\}.
\label{eq:q1-metrics-hours}
\end{align}
Here $P_{\mathrm{peak}}^m$ is reported in kW. The indicator
$\mathbf{1}\{\cdot\}$ equals one when its condition holds and zero otherwise.
Since each interval lasts one hour, $H^m$ is the number of hours above the
cap, reported in h. The tolerance $\varepsilon=2\times10^{-6}$ kWh
matches the numerical tolerance used in the implementation.

Smaller values of $V^m$, $P_{\mathrm{peak}}^m$ and $H^m$ indicate less
residual exceedance, a lower peak and fewer hours above the cap,
respectively. Their improvements are assessed against $\mathcal R$
under the same cap profile. These metrics quantify the import-limitation
service, rather than the DSO's complete monetary net benefit.

The aggregate exceedance $V^{\mathcal I}$ must be distinguished from
the individual contracts' charged exceedance
$\sum_{i,t}x_{i,t}^{\mathcal I}$. The latter can be larger because a
member can exceed its allocated cap while other members leave their
shares unused. In $\mathcal C$, $V^{\mathcal C}=\sum_t x_t^{\mathcal C}$
when the above-cap price strictly exceeds the within-cap price, as in the
base case. For low-penalty sweep points, compute $V^{\mathcal C}$ from
total metered imports using the same definition as for the other models.

All comparisons retain the same member data and battery terminal condition,
and $\mathcal I$ and $\mathcal C$ use identical discount, penalty and
community cap parameters. If the contractual cap is varied, $V^m$ and
$H^m$ are evaluated against the cap of that configuration for all three
models. The actual peak is reported alongside them, so that compliance
with a looser cap is not interpreted as a reduction in peak import.
